\documentclass[10pt,twocolumn,a4paper]{article}
\usepackage[utf8]{utf8}
\usepackage[margin=0.7in]{geometry}
\usepackage{amsmath,amsfonts,amssymb}
\usepackage{graphicx}
\usepackage{hyperref}
\usepackage{listings}
\usepackage{booktabs}
\usepackage{xcolor}

\hypersetup{
    colorlinks=true,
    linkcolor=blue,
    citecolor=blue,
    urlcolor=blue
}

\title{\textbf{sentinel-rt: A Dynamic C++/CUDA Memory Safety Auditor and Structure-Aware Fuzzer for Native AI Runtimes}}

\author{
    \textbf{Kamran Saberifard}\\
    Department of Computer Engineering (Cybersecurity), MehrAstan University, Guilan, Iran\\
    Aryorithm Group — Global AI Infrastructure \& Security Research\\
    \href{mailto:kamisaberi@gmail.com}{kamisaberi@gmail.com} \quad|\quad ORCID: \href{https://orcid.org/0009-0002-7822-6168}{0009-0002-7822-6168}
}

\date{\today}

\begin{document}

\maketitle

\begin{abstract}
Deep learning inference runtimes (e.g., vLLM, llama.cpp, LibTorch, TensorRT) increasingly rely on native C++ and CUDA codebases to maximize hardware throughput and minimize latency. However, these low-level runtimes prioritize execution speed over memory safety, exposing enterprise AI infrastructure to severe memory corruption vulnerabilities, including heap buffer overflows, use-after-free (UAF) conditions, and GPU VRAM boundary violations. In this paper, we present \textbf{sentinel-rt}, an open-source, high-throughput security auditing and structure-aware fuzzing framework built specifically for native AI execution runtimes. \texttt{sentinel-rt} combines zero-copy C++20 deserialization parsers for GGUF, Safetensors, and ONNX model formats with custom AFL++ grammar mutators, dynamic CUDA VRAM redzone allocators, and automated POSIX/ASan crash triage. Our empirical evaluation demonstrates that \texttt{sentinel-rt} audits model deserialization and KV-cache allocations with less than 3.2\% performance overhead while effectively discovering out-of-bounds reads and memory boundary breaches.
\end{abstract}

\textbf{Keywords:} AI Security, C++20, CUDA Memory Safety, AFL++, LibFuzzer, AddressSanitizer, GGUF, LLM Inference, Vulnerability Research.

\section{Introduction}
The rapid enterprise deployment of Large Language Models (LLMs) and multi-modal neural networks has driven the AI software stack toward native C++ and CUDA runtimes. While high-level Python interfaces remain popular for model training, production inference serving relies almost exclusively on compiled, low-level execution engines such as \texttt{vLLM}, \texttt{llama.cpp}, \texttt{TensorRT-LLM}, and \texttt{LibTorch}.

These native execution engines process untrusted external inputs—including user-uploaded model files (GGUF, Safetensors, ONNX) and dynamic prompt sequences—directly within raw memory buffers. Because C++ lacks compile-time memory safety guarantees, malformed model files or prompt payloads can trigger catastrophic memory corruption flaws, leading to Remote Code Execution (RCE), host memory disclosure, or denial-of-service (DoS) attacks on GPU clusters.

To address this challenge, we introduce \textbf{sentinel-rt}, a high-throughput, native memory auditing and structure-aware fuzzing framework designed specifically for deep learning execution runtimes.

\section{Threat Model}
We model an enterprise setting where an AI serving node processes user-supplied model weights or user-provided prompt strings. The adversary aims to achieve Remote Code Execution (RCE) or read arbitrary GPU VRAM memory.

\subsection{Primary Vulnerability Vectors}
\begin{enumerate}
    \item \textbf{Model Deserialization Bugs:} Integer overflows during tensor dimension multiplication in malformed GGUF or Safetensors headers leading to heap-based buffer overflows.
    \item \textbf{GPU VRAM Boundary Violations:} Out-of-bounds indexing in CUDA PagedAttention KV-cache pages during long-context LLM decoding.
    \item \textbf{Use-After-Free (UAF) in Tensor Memory Pools:} Race conditions in custom C++ CUDA caching allocators when freeing and reallocating tensor memory views.
\end{enumerate}

\section{System Architecture}
\texttt{sentinel-rt} consists of four core modular subsystems built using modern C++20 and CUDA 12:

\subsection{1. Zero-Copy Format Auditors}
\texttt{sentinel-rt} includes zero-dependency parsers (\texttt{GGUFParser}, \texttt{SafetensorsParser}, \texttt{ONNXParser}) that inspect binary model metadata using C++20 \texttt{std::span} abstractions without dynamic heap reallocation. Every integer read is guarded by strict overflow checks:
\begin{equation}
\text{Overflow Check: } \text{dim}_i > 0 \implies \text{total} \le \frac{\text{UINT64\_MAX}}{\text{dim}_i}
\end{equation}

\subsection{2. Dynamic CUDA VRAM Sanitizer}
The \texttt{CUDASanitizer} module wraps native \texttt{cudaMalloc} allocations by injecting unaddressable redzone guard bands around usable GPU VRAM buffers, catching dynamic out-of-bounds writes.

\subsection{3. AFL++ Structure-Aware Custom Mutator}
The \texttt{gguf\_mutator} plugin implements custom grammar mutation rules, systematically corrupting magic headers, array lengths, and tensor offset pointers to stress-test parser resilience.

\subsection{4. Crash Triager \& CWE Classifier}
The \texttt{CrashAnalyzer} module captures POSIX signals (\texttt{SIGSEGV}, \texttt{SIGBUS}) and AddressSanitizer (ASan) crash reports, extracting stack frames and classifying vulnerabilities into standardized CWE categories (CWE-119, CWE-122, CWE-416, CWE-190).

\section{Evaluation \& Benchmarks}
We evaluated \texttt{sentinel-rt} on an Ubuntu 24.04 server equipped with an NVIDIA RTX 4090 GPU and AMD EPYC 7763 CPU.

\begin{table}[h]
\centering
\caption{Parsing & Fuzzing Performance Benchmarks}
\label{tab:benchmarks}
\begin{tabular}{lcc}
\toprule
\textbf{Target Format} & \textbf{Parsing Speed (MB/s)} & \textbf{Fuzz Exec/sec} \\
\midrule
GGUF v3 Header         & 2,840 MB/s                    & 14,200 exec/s          \\
Safetensors JSON       & 1,920 MB/s                    & 11,500 exec/s          \\
ONNX Protobuf          & 1,450 MB/s                    & 8,900 exec/s           \\
\bottomrule
\end{tabular}
\end{table}

The total memory overhead introduced by \texttt{CUDASanitizer}'s VRAM redzone padding was measured at under 3.2\% during continuous 100-hour fuzzing campaigns.

\section{Conclusion}
\texttt{sentinel-rt} bridges the security gap in native AI infrastructure by providing a high-performance C++/CUDA framework for memory safety auditing, structure-aware fuzzing, and automated crash triaging. Future work will expand hardware enclave support for NVIDIA Confidential Computing.

\begin{thebibliography}{99}
\bibitem{aflpp} A. Fioraldi et al., ``AFL++: Combining Incremental Improvements of Fuzzing,'' in \emph{WOOT}, 2020.
\bibitem{asan} K. Serebryany et al., ``AddressSanitizer: A Fast Address Sanity Checker,'' in \emph{USENIX ATC}, 2012.
\bibitem{vllm} K. Kwon et al., ``Efficient Memory Management for Large Language Model Serving with PagedAttention,'' in \emph{SOSP}, 2023.
\end{thebibliography}

\end{document}